Now that we have a notion of a euclidean domains, we proceed to establish arithmetic on such rings, notably divisibility and greatest common divisors.

\begin{definition}[Divide, Greatest Common Divisor] \leavevmode \\
    Let $R$ be a commutative ring and let $a,b\in R$ with $b\neq 0$. The element $b$ is said to divide $a$ if
    $$\exists r \in R \st br = a$$
    If $b$ divides $a$, we write $b \mid a$. \\
    The greatest common divisor (GCD) of $a$ and $b$ is a nonzero element of $R$, $d$, such that 
    \begin{enumerate}
        \item $d \mid a$ and $d \mid b$
        \item If $d' \mid a$ and $d' \mid b$, then $d' \mid d$
    \end{enumerate}
    The notation remains the same:
    $$d=\gcd(a,b) \text{ or } d=(a,b)$$
\end{definition}
In the language of ideals if $I$ is the ideal generated by $(a,b)$ then $d$ is a $\gcd(a,b)$ if 
\begin{enumerate}
    \item $I \subseteq (d)$
    \item If $I \subseteq (d')$ then $(d) \subseteq (d')$
\end{enumerate}

\begin{proposition}
    If $a$ and $b$ are nonzero elements of a commutative ring $R$ such that the ideal generated by $a$ and $b$ is a principal ideal $(d)$, then $d$ is a greatest common divisor of $a$ and $b$.
\end{proposition}

\begin{proof}
    Suppose $I = (a,b) = (d)$. Then
    $$a,b\in (d) \implies d\mid a \text{ and } d \mid b$$
    Let $d'\in R \st (a,b) \subseteq (d')$. Then $a \in (d')$ and $b\in (d')$ implies $d' \mid a$ and $d' \mid b$. Since $(d) = (a,b),$ $(d) \subseteq (d')$ which implies $d\in (d')$. More specifically, $d' \mid d$. Hence, $d$ is a greatest common divisor of $a$ and $b$.
    \qed
\end{proof}

Just like for standard integers, the greatest common divisor is unique.

\begin{proposition}
    Let $R$ be an integral domain. If $d$ and $d'$ in $R$ both generate the same principal ideal, i.e. $(d) = (d')$, then $d=ud'$ where $u$ is a unit in $R$. In particular, if $d$ and $d'$ are both greatest common divisors of $a$ and $b$, then $d' = wd$ where $w$ is some unit in $R$.
\end{proposition}

\begin{proof}
    This is trivial if $d$ or $d'$ is zero. Assume $d$ and $d'$ are both nonzero. Since $(d) = (d')$, we have that $d' \in (d)$ and so $\exists x \in R \st dx = d'$. Since $d \in (d')$, $\exists y \in R \st d'y = d$. This gives
    \begin{align*}
        d'yx = d' &\implies d'yx - d' = 0 \\
        &\implies d'(yx - 1) = 0
    \end{align*}
    Since $d' \neq 0$ and $R$ is an integral domain, we know that $yx - 1 = 0 \implies yx = 1$. Hence, both $x$ and $y$ are units.
    \qed
\end{proof}